\begin{frame}{16개 층을 쌓아도: 풀링이 시야를 2배로}
  \begin{columns}[T]
    \begin{column}{0.56\textwidth}
      \centering
      \includegraphics[width=\linewidth]{ch10_2_receptive_field.png}
      \\[0.3em]{\footnotesize 16개 층을 쌓아도 풀링 없이
      스트라이드 1 합성곱만 쌓으면 수용장이 33px(입력 32px과
      비슷)에 그치는 반면, [Conv3, Conv3, Pool2] 블록 4개를
      쌓으면 \textbf{76px} --- 이미지의 2배. 스트라이드 2 풀링
      한 번이 그 뒤의 모든 확장을 2배 스케일업.}
    \end{column}
    \begin{column}{0.44\textwidth}
      \begin{itemize}
        \item Week 4.2 MLP에서 ``은닉층을 쌓을수록 결정경계가
          더 복잡한 형태''라는 관측과 \textbf{같은 원리}
        \item \textbf{얕은 층}은 지역적 패턴(선, 모서리)을,
          \textbf{깊은 층}은 그 지역적 패턴을 \textbf{넓은
          시야에서 조합}한 고수준 특징(눈, 바퀴)을 감지
      \end{itemize}
    \end{column}
  \end{columns}
\end{frame}
